%-------------------------------------------------------
%	CLASS, PACKAGES AND OTHER DOCUMENT CONFIGURATIONS
%-------------------------------------------------------
\documentclass{article}[
	a4paper, % Paper size
	12pt,    % Font size
	%unnumberedsections, % Uncomment for no section numbering
]
%-------------------------------------------------------
%	CLASS CONFIGURATION
%-------------------------------------------------------
\usepackage{etoolbox} % Required for conditional logic and easily changing commands

\newtoggle{unnumberedsections} 			% Create toggle for a class option
\settoggle{unnumberedsections}{false} 	% Default value for the class option
\DeclareOption{unnumberedsections}{\settoggle{unnumberedsections}{true}} % Set the class option toggle if the class option was used in the template

\DeclareOption*{\PassOptionsToClass{\CurrentOption}{article}} 			 % Pass through any extra options specified to the base class
\ProcessOptions\relax 					% Process class options

%-------------------------------------------------------
%	REQUIRED PACKAGES AND MISC CONFIGURATIONS
%-------------------------------------------------------
\usepackage{graphicx} 	     % Required for including images
\graphicspath{{figures/}{./}} % Figures path

\setlength{\parindent}{0pt} % Paragraph indentation
\setlength{\parskip}{10pt}  % Vertical space between paragraphs

\usepackage[bottom, hang]{footmisc} % Force footnotes to the bottom of the page and to the left margin
\setlength{\footnotemargin}{6pt}    % Horizontal space between the footnote marker and text

\usepackage[document]{ragged2e} % Left align all text in the document (i.e. have a ragged right margin)
\usepackage{lastpage} 			% Required to determine the total number of pages

\AtBeginEnvironment{quote}{\itshape\rmfamily} % Style quotations

\usepackage{appendix}  % Required for the appendices section
\usepackage{float} 		% Allows more precisely positioning floats e.g. \begin{figure}[H]
	

%-------------------------------------------------------
%	MARGINS
%-------------------------------------------------------
\usepackage[
	top=3cm, 	  			% Top margin
	bottom=3cm,  			% Bottom margin
	left=2cm,	  			% Left margin
	right=7.25cm, 			% Right margin
	marginparwidth=4.25cm,  % Width of the large right margin (sidebar) content
	marginparsep=1cm, 		% Padding between the main text and the large right margin (sidebar)
	footskip=1.5cm, 		% Space from the bottom margin to the baseline of the footer
	headsep=0.8cm, 			% Space from the top margin to the baseline of the header
	headheight=1cm, 		% Height of the header
	%showframe % Debugging frames, uncomment to show frames around margins
]{geometry}
\usepackage[oneside]{sidenotes} % Required for outputting content to the margin\

% Restyle \sidenote command from sidenotes package to reduce the font size
\makeatletter
\RenewDocumentCommand\sidenotetext{o o +m}{%      
	\IfNoValueOrEmptyTF{#1}{%
		\@sidenotes@placemarginal{#2}{\textsuperscript{\thesidenote}\footnotesize#3}%
		\refstepcounter{sidenote}%
	}{%
		\@sidenotes@placemarginal{#2}{\textsuperscript{#1}\footnotesize#3}%
	}%
}
\makeatother
% Environment to output to the full page width (excluding page margins)
\newenvironment{fullwidth}{ 
	% Pull into the right margin by the margin width and margin skip width
	\begin{adjustwidth}{}{\dimexpr-\marginparwidth-\marginparsep\relax} 
}{
	\end{adjustwidth}
}
% Environment to output content at two-thirds of the whole page width
\newenvironment{twothirdswidth}{ 
	% Pull into the right margin by half the margin width and half the margin skip width
	\begin{adjustwidth}{}{\dimexpr-0.5\marginparwidth-0.5\marginparsep\relax} 
}{
	\end{adjustwidth}
}

% Calculate the full text width (including the large right margin and margin padding)
\newlength{\fulltextwidth}
\setlength{\fulltextwidth}{\dimexpr\textwidth+\marginparwidth+\marginparsep\relax}
% New command to output an unnumbered sidenote, the optional argument is a vertical offset and the required parameter is the sidenote text
\newcommand{\nonumsidenote}[2][]{\sidenote[\hspace{0pt}][#1]{#2}} 
% New command to output a quotation sidenote, the optional argument is a vertical offset and the required parameter is the quotation text
\newcommand{\sidenotequote}[2][]{\sidenote[\hspace{0pt}][#1]{\rmfamily\itshape#2}} 
%-------------------------------------------------------
%	FONTS
%-------------------------------------------------------
\usepackage[utf8]{inputenc} 	  % Required for inputting international characters
\usepackage[T1]{fontenc} 		  % Output font encoding for international characters
\usepackage{mathpazo} 			  % Use Palatino for math
\usepackage[sfdefault]{plex-sans} % Use IBM Plex Sans as the sans font and make it the default
\usepackage{plex-serif} 		  % Use IBM Plex Serif as the serif font
\usepackage{plex-mono} 			  % Use IBM Plex Mono as the mono font

% ----------Rule Definations---------- {el}{l}{sb}
\newcommand{\textel}[1]{{\fontseries{el}\selectfont #1}} % ExtraLight weight command
\newcommand{\textl}[1]{{\fontseries{l}\selectfont #1}}   % Light weight
\newcommand{\textsb}[1]{{\fontseries{sb}\selectfont #1}} % SemiBold weight

%-------------------------------------------------------
%	HEADERS AND FOOTERS
%-------------------------------------------------------
\usepackage{fancyhdr} % Required for customizing headers and footers
\usepackage{multicol} % Used for multipule colums
\pagestyle{fancy} 	  % Enable custom headers and footers

\renewcommand{\headrulewidth}{0.5pt} % Top horizontal rule thickness

\fancyhf{} 			  % Clear default headers/footers

\fancyhead[L]{\reporttitle} 	   % Left header
\fancyhead[C]{} 				   % Center header
\fancyhead[R]{\rightheadercontent} % Right header

\fancyfoot[L]{} 							  % Left footer
\fancyfoot[C]{\thepage~of \pageref{LastPage}} % Center footer
\fancyfoot[R]{} 							  % Right footer

\usepackage{calc} 	  % Allows arithmetic in LaTeX commands
\newlength{\marginoffset}
\setlength{\marginoffset}{\marginparwidth + \marginparsep} % Calculate width of margin and buffer space before margin
\fancyheadoffset[R]{\marginoffset} 						   % Offset the header to account for the margin width
\fancyfootoffset[R]{\marginoffset} 						   % Offset the footer to account for the margin width

%-------------------------------------------------------
%	SECTIONS
%-------------------------------------------------------
\usepackage{titlesec} 		    % Required for modifying sections
\iftoggle{unnumberedsections}{  % Conditional logic for the unnumbered sections class options
	\setcounter{secnumdepth}{0} % Don't number sections at any level
}{
	\setcounter{secnumdepth}{3} % Number sections down to subsubsections
}
\titleformat
	{\section} % Section type being modified
	[block]    % Section layout type, can be: hang, block, display, runin, leftmargin, rightmargin, drop, wrap, frame
	{\raggedright\LARGE\bfseries} % Text formatting of the whole section, i.e. label and title
	{\thesection} % Section label (e.g. number) and its formatting
	{5pt} 		  % Horizontal space between the section label and title
	{} % Code before the section title
	[] % Code after the section title

\titlespacing*{\section}{0pt}{\baselineskip}{4pt} % Spacing around section titles, the order is: left, before and after

%------------------------------------------------

\titleformat
	{\subsection} % Section type being modified
	[block] 	  % Section layout type, can be: hang, block, display, runin, leftmargin, rightmargin, drop, wrap, frame
	{\raggedright\large\bfseries} % Text formatting of the whole section, i.e. label and title
	{\thesubsection} % Section label (e.g. number) and its formatting
	{5pt} 			 % Horizontal space between the section label and title
	{} % Code before the section title
	[] % Code after the section title

\titlespacing*{\subsection}{0pt}{0.75\baselineskip}{4pt} % Spacing around section titles, the order is: left, before and after

%------------------------------------------------

\titleformat
	{\subsubsection} % Section type being modified
	[block] 		 % Section layout type, can be: hang, block, display, runin, leftmargin, rightmargin, drop, wrap, frame
	{\raggedright\bfseries\normalsize} % Text formatting of the whole section, i.e. label and title
	{\thesubsubsection} % Section label (e.g. number) and its formatting
	{5pt} 				% Horizontal space between the section label and title
	{} % Code before the section title
	[] % Code after the section title

\titlespacing*{\subsubsection}{0pt}{0.5\baselineskip}{0pt} % Spacing around section titles, the order is: left, before and after

%------------------------------------------------

\titleformat
	{\paragraph} % Section type being modified
	[runin] 	 % Section layout type, can be: hang, block, display, runin, leftmargin, rightmargin, drop, wrap, frame
	{\bfseries}  % Text formatting of the whole section, i.e. label and title
	{} 			 % Section label (e.g. number) and its formatting
	{0pt} 		 % Horizontal space between the section label and title
	{} 			 % Code before the section title
	[~~~~] 		 % Code after the section title

\titlespacing*{\paragraph}{0pt}{0pt}{0pt} % Spacing around section titles, the order is: left, before and after

%-------------------------------------------------------
%	TABLE OF CONTENTS
%-------------------------------------------------------
\usepackage{titletoc}    % Required for manipulating the table of contents
\setcounter{tocdepth}{2} % Show entries in the table of contents down to subsections

\renewcommand{\contentsname}{Table of Contents} % Rename the table of contents section heading

\titlecontents{section} % Section type being modified
	[0pt] 				% Left indentation
	{\vspace{-2pt}} 	% Before code
	{\textsb{\thecontentslabel}~~\textsb} % Formatting of numbered sections of this type
	{\textsb} 			% Formatting of numberless sections of this type
	{\titlerule*[6pt]{.}\textsb{\contentspage}} % Formatting of the filler to the right of the heading and the page number
	[\vspace{2pt}] 		% After code

%------------------------------------------------

\titlecontents{subsection} % Section type being modified
	[0.5cm] 			   % Left indentation
	{\vspace{-8pt}} 	   % Before code
	{\textsb{\thecontentslabel}~~} % Formatting of numbered sections of this type
	{} 					   % Formatting of numberless sections of this type
	{\titlerule*[6pt]{.}\contentspage} % Formatting of the filler to the right of the heading and the page number
	[] 					   % After code
	
%------------------------------------------------

\titlecontents{subsubsection} % Section type being modified
	[1.35cm] 				  % Left indentation
	{\vspace{-4pt}} 		  % Before code
	{\textsb{\thecontentslabel}~~} % Formatting of numbered sections of this type
	{} 						  % Formatting of numberless sections of this type
	{\titlerule*[6pt]{.}\contentspage} % Formatting of the filler to the right of the heading and the page number
	[] 						  % After code

%-------------------------------------------------------
%	COLORS
%-------------------------------------------------------
\usepackage[usenames, svgnames]{xcolor,colortbl} % Required for defining and using custom colors
%\definecolor{cyan}{RGB}{0, 255, 255} % Define custom colors in this format

%-------------------------------------------------------
%	TABLES
%-------------------------------------------------------
\usepackage{booktabs}   % Required for better horizontal rules in tables (\toprule/bottomrule/addlinespace)
\usepackage{rotating}   % Rotating table labels
\usepackage{array} 	    % Required for manipulating table columns
\usepackage{tabularray} % LaTeX3, used for colspec / rowspec
\renewcommand{\arraystretch}{1.25} % Increase the height of table rows

% ----------Rule Definations----------
\newcolumntype{R}[1]{>{\raggedleft\arraybackslash}p{#1}}  % New right-aligned paragraph column type
\newcolumntype{L}[1]{>{\raggedright\arraybackslash}p{#1}} % New left-aligned (no justification) paragraph column type
\newcolumntype{C}[1]{>{\centering\arraybackslash}p{#1}}   % New centered paragraph column type

%-------------------------------------------------------
%	CAPTIONS
%-------------------------------------------------------
\usepackage{caption} 		   % Required for customizing captions
\captionsetup{skip=6pt} 	   % Vertical whitespace between figures/tables and the caption (default is 10pt)
\captionsetup{font={bf,small}} % Define caption font style
%\captionsetup{justification=raggedright, singlelinecheck=false} % Left-align captions

% The sidenotes package defines 4 custom caption styles for wide tables, wide figures, margin tables and margin figures
% To style these the same as in-text tables and figures, set these 4 styles to the same styling values
\DeclareCaptionStyle{widetable}{font={bf, footnotesize}, justification=raggedright, singlelinecheck=false}
\DeclareCaptionStyle{widefigure}{font={bf, footnotesize}, justification=raggedright, singlelinecheck=false}
\DeclareCaptionStyle{marginfigure}{font={bf, scriptsize}, justification=raggedright, singlelinecheck=false}
\DeclareCaptionStyle{margintable}{font={bf, scriptsize}, justification=raggedright, singlelinecheck=false}

%-------------------------------------------------------
%	LISTS
%-------------------------------------------------------
\usepackage{enumitem} 			  % Required for list customization

\setlist{itemsep=0pt, topsep=0pt} % Customize spacing around and inside lists

\renewcommand{\labelenumi}{\arabic{enumi}.}    % Customize the first level numbered list characters
\renewcommand{\labelenumii}{\alph{enumii}.}    % Customize the second level numbered list characters
\renewcommand{\labelenumiii}{\roman{enumiii}.} % Customize the third level numbered list characters

\renewcommand{\labelitemi}{\Large\raisebox{-1.5pt}{\textbullet}} % Customize the first level bullet point character
\renewcommand{\labelitemii}{\Large\raisebox{-1.5pt}{--}} 		 % Customize the second level bullet point character
\renewcommand{\labelitemiii}{\Large\raisebox{-1.5pt}{-}} 		 % Customize the third level bullet point character

%-------------------------------------------------------
%	LINKS
%-------------------------------------------------------
\usepackage{hyperref} % Required for links

\hypersetup{
	colorlinks=true, % Whether to color the text of links
	urlcolor=black,  % Color for \url and \href links
	linkcolor=black, % Color for \nameref links
	citecolor=black, % Color of reference citations
}

%-------------------------------------------------------
%	CODE LISTINGS
%-------------------------------------------------------
\usepackage{listings} % Required for formatted code blocks

\lstset{
	numbers=left, 		  		% Line numbers location
	numbersep=6pt,		  		% Whitespace to the right of line numbers
	numberstyle=\tiny\it, 		% Line numbers font style
	frame=single, 		  		% Border
	basicstyle=\tt\scriptsize,  % Code font style
	breaklines=true,      		% Break long lines into multiple lines automatically
	aboveskip=1.5\baselineskip, % Vertical whitespace before the listing
	tabsize=3, 			  		% How many space to a tab
}

%-------------------------------------------------------
%	CUSTOM COMMANDS
%-------------------------------------------------------
\newcommand{\reporttitle}[1]{\renewcommand{\reporttitle}{#1}}
\newcommand{\reportsubtitle}[1]{\renewcommand{\reportsubtitle}{#1}}
\newcommand{\reportauthors}[1]{\renewcommand{\reportauthors}{#1}}
\newcommand{\reportdate}[1]{\renewcommand{\reportdate}{#1}}
\newcommand{\rightheadercontent}[1]{\renewcommand{\rightheadercontent}{#1}}

%-------------------------------------------------------
%	BIBLIOGRAPHY
%-------------------------------------------------------
\usepackage[
	backend=biber,     % Use the biber backend for compiling the bibliography
	citestyle=numeric, % In-text citation style
	bibstyle=numeric,  % Bibliography style
]{biblatex}

\newcommand{\sidecite}[2][0pt]{\sidenote[][#1]{\textcite{#2}}}     % Parenthetical citations in the margin
\newcommand{\fullsidecite}[2][0pt]{\sidenote[][#1]{\fullcite{#2}}} % Full citations in the margin

\addbibresource{sample.bib} % BibLaTeX bibliography file

%------------------------------------------------
%	REPORT INFORMATION
%------------------------------------------------
\reporttitle{CompTIA Security+} 			    % The report title to appear on the title page and page headers
\reportsubtitle{Personal Notes {\color{red} Exam Objectives}} % Report subtitle, include new lines if needed
\reportauthors{\textit{Author:}\\\smallskip Scott \textit{Skrobel}} 		% Report authors/group/department, include new lines if needed
\reportdate{\today}														    % Report date, include new lines for additional information if needed
\rightheadercontent{\includegraphics[width=3cm]{logo.png}} 					% The content in the right header, leave empty for no right header content

%%%%%%%%%%%%%%%%%%%%%%%%%%%%%%%%%%%%%%%%%%%%%%%%%%%%%%%%%%%%%%%%
%							START DOC
%%%%%%%%%%%%%%%%%%%%%%%%%%%%%%%%%%%%%%%%%%%%%%%%%%%%%%%%%%%%%%%%
\begin{document}
%------------------------------------------------
%	Title Page
%------------------------------------------------
\thispagestyle{empty} % Suppress headers and footers on this page

\begin{fullwidth}     % Use the whole page width
	\vspace*{-0.075\textheight} 	           % Pull logo into the top margin
	%\hfill\includegraphics[width=5cm]{logo.png} % Company logo
	\vspace{0.15\textheight} 				   % Vertical whitespace
	% Report title, intentionally at less than full width for nice wrapping. Adjust width of \parbox and the font size as needed for your title to look good.
	\parbox{1.9\fulltextwidth}{\fontsize{50pt}{52pt}\selectfont\raggedright\textbf{\reporttitle}\par} 
	
	\vspace{0.03\textheight} 					  % Vertical whitespace	
	{\LARGE\textit{\textbf{\reportsubtitle}}\par} % Subtitle
	\vspace{0.04\textheight} 					  % Vertical whitespace
	{\color{black} \rule{\linewidth}{0.5mm}}      % Color Horizontal Line
	\vfill 										   % Vertical whitespace
	
	{\Large\reportauthors\par} % Report authors, group or department

	\vfill\vfill\vfill 			  % Vertical whitespace
	\begin{flushright}
		
	{\large\reportdate\par}    % Report date
	\end{flushright}
\end{fullwidth}

\newpage
%------------------------------------------------
%	TABLE OF CONTENTS
%------------------------------------------------

\begin{twothirdswidth} % Content in this environment to be at two-thirds of the whole page width
	\tableofcontents   % Output the table of contents, automatically generated from the section commands used in the document
\end{twothirdswidth}

\newpage

%------------------------------------------------
%	SECTION 1: SUBSECTIONS: 1-8
%------------------------------------------------
\section{Threats, Attacks and Vulnerabilities}
\subsection{Compare and contrast types of attacks}
\begin{fullwidth}
    \begin{description}\itemsep2pt
        \item[Phishing]: Social engineering tactic to acquire personal information from a fake email with a clickable link.
        \item[Smishing]: (SMS phishing) The use of deceptive text messages into divulging sensitive information.
        \item[Vishing]: (Voice Phishing) Impersonates a trusted entity, such as a bank to trick into giving information.
        \item[Spam]: Unsolicited inappropriate messages sent with the purpose of spreading malware, advertising, phishing. 
        \item[Spear phishing]: Targeted type of phishing attack to make the scam convincing, often with insider information.
        \item[Dumpster diving]: Searching through an organization's or individual's trash to find sensitive information.
        \item[Shoulder surfing]: Observing a victim’s screen or keyboard to obtain sensitive information.
        \item[Pharming]: Manipulating the DNS system to resolve fake domain names, to lead them to a fake website.
        \item[Tailgating]: Physical security breach by following authorized person to get access to secure areas.
        \item[Eliciting information]: Psychological tactics to encourage individuals to share their knowledge willingly. 
        \item[Whaling]: Spear phishing for high-profile executives in an organization.
        \item[Prepending]: Organizing, manipulating, and structuring data in various applications.
        \item[Identity fraud]: An individual wrongfully obtains and uses someone else's personal data in a deceptive manner.
        \item[Invoice scams]: Impersonation of a legit business to deceive individuals into paying fraudulent invoices. 
        \item[Credential harvesting]: Tricking someone into disclosing login credentials to access sensitive info.
        \item[Reconnaissance]: Initial phase to gather intelligence via passive and active techniques.
        \item[Hoax]: Fabrication intended to deceive or trick individuals into believing false information or events.
        \item[Impersonation]: Masquerading as a legitimate user or entity to gain unauthorized access to information.
        \item[Watering hole attack]: Infecting a commonly visited website of a targeted specific group.
        \item[Typosquatting]: URL hijacking or domain squatting of fake domains which resemble a legit website.
        \item[Pretexting]: Fabricated scenario involving direct interaction to obtain sensitive information.
        \item[Influence campaigns]: Coordinated effort to shape public opinion, influence perceptions, and manipulate.
        \begin{itemize}\itemsep2pt
            \item[]\textbf{Hybrid warfare}: Blends conventional warfare tactics with unconventional method
        \end{itemize}
        \item[Principles] (reasons for effectiveness)
        \begin{itemize}\itemsep2pt
            \item \textbf{Authority}: The actor acts as an individual of authority
            \item \textbf{Intimidation}: Frightening or threatening the victim.
            \item \textbf{Consensus}: Convince based on what’s normally expected.
            \item \textbf{Scarcity}: Limited resources and time to act.
            \item \textbf{Familiarity}: The victim is well known.
            \item \textbf{Trust}: Gain their confidence, be their friend.
            \item \textbf{Urgency}: Limited time to act, rush the victim.
        \end{itemize}
    \end{description}
\end{fullwidth}

\newpage
\subsection{Given a scenario, analyze indicators of compromise}
\subsubsection*{\color{red}Malware}
\begin{fullwidth}
    \begin{description}\itemsep2pt
        \item[Ransomware:] Denies access to a computer system or data until a ransom is paid.
        \item[Trojan:] A form of malware that pretends to be a harmless application.
        \item[Worm:] A self-contained infection that can spread itself through networks, emails, and messages.
        \item[PUP's:] Potentially Unwanted Programs software applications that may exhibit undesirable characteristics.
        \item[Memory-resident malware:] Operates primarily in a computer's volatile memory (RAM) rather than with files
        \item[Command and control:] (C2) Centralized server used by attackers to manager compromised devices.
        \item[Bots:] AI inside an infected machine performs specific actions as a part of a larger entity known as a botnet.
        \item[Cryptomalware:] A malicious program that encrypts programs and files on the computer to extort money.
        \item[Logic Bomb:] A malicious program that lies dormant until a specific date or event occurs.
        \item[Spyware:] Software that installs itself to spy and sends stolen info back to the host machine.
        \item[Keyloggers:] A malicious program that saves all of the keystrokes of the infected machine.
        \item[Remote Access Trojan] (RAT) A remotely operated Trojan.
        \item[Rootkit:] A backdoor program that allows full remote access to a system.
        \item[Backdoor:] Allows for full access to a system remotely.
    \end{description}
\end{fullwidth}

\subsubsection*{\color{blue}Password Attacks}
\begin{fullwidth}
    \begin{description}\itemsep2pt
        \item[Brute Force:] Systematically trying a large number of possible passwords.
        \begin{itemize}\itemsep2pt
            \item \textbf{Offline}: Attempting to crack a password hash without directly interacting with the target system, rather on their own independant computer.
            \item \textbf{Online}: Attempting to guess a user's password by repeatedly trying different combinations.
        \end{itemize}
        \item[Sprying:] A Type of brute-force attack by attempting to authenticate with commonly used passwords. Small number of passwords against many accounts.
        \item[Dictionary:] A password attack that creates encrypted versions of common dictionary words and then compares them against those in a stolen password file. Guessing using a list of possible passwords. 
        \item[Rainbow Table:] Large pregenerated data sets of encrypted passwords used in password attacks.
        \item[Plaintext/Unencrypted:] The attacker has both the plaintext and its encrypted version.
    \end{description}
\end{fullwidth}

\subsubsection*{\color{purple}Physical attacks}
\begin{fullwidth}
    \begin{description}
        \item[Malicious Flash Drive:] A storage device loaded with malware.
        \item[Serial Bus (USB) cable:] A USB cable designed to compromise systems upon connection.
        \item[Card cloning]: Creating a copy of a credit or other card with stolen data.
        \item[Skimming:] Stealthily capturing and storing all the details stored on your card’s magnetic stripe.
    \end{description}
\end{fullwidth}

\subsubsection*{\color{purple}Adversarial AI}
\begin{fullwidth}
    \begin{description}
        \item[Tainted Training Data for ML:] Modifying the data used to train machine learning models to cause misclassifications or errors.
        \item[Security of Machine Learning Algorithms:] Ensuring ML algorithms are protected against manipulation and attacks.
    \end{description}
\end{fullwidth}

\subsubsection*{\color{purple}Other}
\begin{fullwidth}
    \begin{description}
        \item[Supply-chain Attacks] Targeting less-secure elements in the supply network to compromise a primary target.
        \item[Cloud-based vs. On-premises Attacks:] Security incidents occurring either in a cloud infrastructure or on locally hosted (on-premises) resources.
    \end{description}
\end{fullwidth}

\subsubsection*{\color{purple}Cryptographic Attacks}
\begin{fullwidth}
    \begin{description}
        \item[Birthday:] Exploiting the probability of two distinct inputs having the same output.
        \item[Collision:] Finding two different inputs that provide the same output.
        \item[Downgrade:] Forcing a system to fall back to a less secure version to exploit vulnerabilities.]
    \end{description}
\end{fullwidth}
\newpage
\begin{fullwidth}    
    \subsection{Given a scenario, analyze potential indicators associated with application attacks}
    \begin{description}
        \item[Privilege Escalation:] An attack that exploits a vulnerability that allows them to gain access to resources that they normally would be restricted from accessing.(imagine logging into a computer as a guest account and having access to admin power)
        \item[Cross-site Scripting:]\textbf{(XXS)} It’s like sneaking a secret note into a bunch of official letters. You insert malicious scripts into websites, which then run on another user’s browser, stealing information or performing actions on their behalf without them knowing.
        \item[Injections:] Occurs when processing invalid data, inserts code into vulnerable program and changes the course of execution.
        \begin{itemize}
            \item \textbf{Structured query language (SQL Injection):} Inserting SQL code into a query to manipulate a database (i.e. to view, edit, or delete data).
            \item \textbf{Dynamic-link library(DLL Injection):} Inserting code into a running process by taking advantage of Dynamic Link Libraries used by software.
            \item \textbf{ Lightweight Directory Access Protocol (LDAP Injection):} Manipulating Lightweight Directory Access Protocol queries (used for organizing/finding user or device data in networks).
            \item \textbf{ Extensible Markup Language (XML Injection):} Inserting elements into an XML document to exploit the structure and logic of an application.
        \end{itemize}
        \item[Pointer/object dereference:] Imagine forgetting to check who’s knocking at the door and just letting them in — failing to validate who or what a pointer is pointing to can allow unauthorized access or crashes.
        \item[Directory Traversal:] It’s like navigating through a building’s restricted areas by exploiting weak security, accessing unauthorized files/folders in a system.
        \item[Buffer Overflows:] Imagine pouring water into a glass until it overflows, only here, excessive data overflows into other memory areas, potentially allowing malicious code execution. 
        \item[Race Conditions \& Time of Check/Time of Use:] Two actions racing to utilize a resource and whoever wins could impact the system. If malicious action wins, it can exploit the time gap between checking a condition and using a resource.
        \item[Error Handling:] How a system responds to unexpected inputs or conditions — poor error handling might expose sensitive information or pathways to attacks.
        \item[Input Handling:] Not checking or sanitizing input properly could allow harmful data into a system, causing malfunctions or unauthorized activities.
        \item[Replay Attack \& Session Replays:] Replaying is resending data (like login credentials) intercepted earlier to gain unauthorized access. Session replays involve capturing and reusing session identifiers, allowing attackers to impersonate legitimate users.
        \item[Integer Overflow:] It’s like an odometer rolling over to zero after reaching its maximum value, only here, exceeding numerical storage capacity might cause erratic system behavior.
        \item[Request Forgeries:] Tricking a user or system into performing actions without knowing:
        \begin{itemize}
            \item \textbf{Server-Side Request Forgery (SSRF):} Making a server unknowingly perform actions on behalf of an attacker.
            \item \textbf{Cross-Site Request Forgery (CSRF):} Making a user’s browser perform an unwanted action on a site where they are authenticated.
        \end{itemize}
        \item[API Attacks:] Exploiting vulnerabilities in APIs — essentially, pathways that let different software components communicate — to interfere with an application’s functionality or steal data.
        \item[Resource Exhaustion:] Draining a system’s resources (like memory or processing power) to slow it down or cause a failure, making it vulnerable to other attacks.
        \item[Memory Leak:] Continually using up memory without releasing it back, like continually filling a basket with apples and never emptying it, which eventually causes slowdowns or crashes.
        \item[Secure Sockets Layer (SSL) stripping:] Downgrading a secure HTTPS connection to an unsecured HTTP connection, making data transmission vulnerable to interception.
        \item[Driver Manipulation:] 
        \begin{itemize}
            \item \textbf{Shimming}: Using extra code (a shim) to make a driver run in environments it’s not compatible with, potentially opening security gaps.
            \item \textbf{Refactoring}: Changing the driver’s internal structure without altering its external behavior, potentially introducing vulnerabilities.
        \end{itemize}
        \item[Pass the Hash:] Using a user’s hash (a type of encrypted password) to authenticate with a service without knowing the actual password.
        
    \end{description}
    
\end{fullwidth}
\newpage
\begin{fullwidth} 
    \subsection{Given a scenario, analyze potential indicators associated with network attacks}
    \subsubsection*{\color{red}Wireless}
    \begin{description}
        \item[Evil Twin:] Imagine someone impersonating your Wi-Fi network to trick devices into connecting to it. It’s an “evil twin” of your legit Wi-Fi, stealing data and spying on users.
        \item[Rogue Access Point:] An unauthorized Wi-Fi access point, maybe added by an employee or attacker, which can bypass security settings.
        \item[Bluesnarfing:] Stealing information from Bluetooth-enabled devices by exploiting vulnerabilities in their Bluetooth connection.
        \item[Bluejacking:] Sending unsolicited messages to a Bluetooth device, mostly harmless but potentially annoying.
        \item[Disassociation:] Interrupting the Wi-Fi connection between a device and a network, causing disruptions.
        \item[Jamming:] Flooding a frequency (like Wi-Fi or cell frequencies) to block communications.
        \item[Radio frequency identification (RFID):] A tech that uses radio waves for tracking and identification but can be exploited to illicitly read information.
        \item[Near-field communication (NFC):] A way to wirelessly share data over short distances, like payment info, which can be exploited for unauthorized data access.
        \item[Initialization Vector (IV):] A random number used in cryptography for preventing predictability in encrypted data, but if not handled properly, can be a vulnerability. 
    \end{description}
    \subsubsection*{\color{Purple}Layer 2 Attacks}
    \begin{description}
        \item[ARP Poisoning:] Confusing network devices by sending fake Address Resolution Protocol messages, redirecting traffic through an attacker’s device.
        \item[Media access control (MAC) Flooding:] Overflowing the network switch with too many Media Access Control addresses, forcing it into acting like a basic hub and revealing internal data traffic.
        \item[MAC Cloning:] Copying a legit MAC address to impersonate a network device.
    \end{description}

    \subsubsection*{\color{blue}Domain Name System (DNS)}
    \begin{description}
        \item[Domain Hijacking:] Taking control of a domain away from the rightful owner, often for malicious activities.
        \item[DNS Poisoning:] Providing false DNS responses to redirect a user’s traffic to malicious sites.
        \item[URL Redirection:] Manipulating URLs to direct users to unintended pages, often for phishing.
        \item[Domain Reputation:] How trustworthy a domain is, based on its past actions and security posture.
    \end{description}

   
    \begin{description}
        \item[Distributed denial-of-service (DDoS):] Overwhelming a target, such as a website, with a flood of internet traffic, making it unavailable to users. Variants include targeting network, application, or operational technology layers.
        \item[On-Path Attack (Man-in-the-Middle):] This is like eavesdropping, where the attacker intercepts and possibly alters the communication between two parties without them knowing.
    \end{description}
    
    \subsubsection*{\color{green}Malicious code or script execution}
    \begin{description}
        \item[PowerShell, Python, Bash:] Different scripting languages that can be used to automate tasks or exploit vulnerabilities.
        \item[Macros, VBA:] Automated scripts, often in Office documents, that can be exploited to run malicious code.
    \end{description}
    
\end{fullwidth}

\newpage
\begin{fullwidth} 
    \subsection{Different threat actors, vectors, and intelligence sources}
    \subsubsection*{\color{red}Actors and Threats}
    \begin{description}
        \item[dvanced Persistent Threat (APT):] Highly skilled attackers, often funded by governments, who aim to stealthily infiltrate and stay in networks for a long time, usually for espionage.
        \item[Insider Threats:] People inside an organization (like employees or contractors) who pose security risks, either maliciously or inadvertently.
        \item[State Actors:] Hackers sponsored by national governments to engage in cyber espionage, warfare, or sabotage.
        \item[Hacktivists:] Individuals or groups hacking for political or social reasons rather than financial gain.
        \item[Script Kiddies:] Inexperienced hackers who use pre-written scripts or tools to perform attacks, without much understanding of how they work.
        \item[Criminal Syndicates:] Organized crime groups engaging in cybercrime for financial gain.
        \item[Hackers:] People who find and exploit vulnerabilities in systems. They can be:
        \begin{itemize}
            \item \textbf{Authorized:} Have permission to access.
            \item \textbf{Unauthorized:} No permission to access.
            \item \textbf{Semi-authorized:} Somewhere in between; maybe they had permission at one point or for certain tasks.
        \end{itemize}
        \item[Shadow IT:] Unauthorized tech solutions used inside an organization without the IT department’s knowledge or approval.
        \item[Competitors:] Business rivals who might engage in cyber tactics to gain a competitive edge.
    \end{description}

    \subsubsection*{\color{red}Attributes of Actors}
    \begin{itemize}
        \item \textbf{Internal/External:} Are they inside or outside the organization?
        \item \textbf{Level of Sophistication:} How skilled are they?
        \item \textbf{Resources:} What tools, money, or people do they have at their disposal?
        \item \textbf{Intent/Motivation:} Why are they doing what they’re doing?
    \end{itemize}

    \subsubsection*{\color{purple}Vectors}
    \begin{itemize}
        \item \textbf{Direct Access:} Physically accessing systems.
        \item \textbf{Wireless:} Via Wi-Fi, Bluetooth, etc.
        \item \textbf{Email:} Think phishing or malware attachments.
        \item \textbf{Supply Chain:} Targeting suppliers or service providers.
        \item \textbf{Social Media:} Spreading malware or misinformation.
        \item \textbf{Removable Media:} USB drives, DVDs, etc.
        \item \textbf{Cloud:} Exploiting vulnerabilities in cloud services.
    \end{itemize}

    \subsubsection*{\color{purple}Threat Intelligence Sources}
    \begin{itemize}
        \item \textbf{Open-Source Intelligence (OSINT):} Publicly available info.
        \item \textbf{Vulnerability Databases:} Listings of known security vulnerabilities.
        \item \textbf{Public/Private Information-Sharing Centers:} Organizations that share threat data.
        \item \textbf{Dark Web:} A part of the internet not indexed by search engines, often hosting illegal activities.
        \item \textbf{Indicators of Compromise:} Signs that a breach has occurred.
        \item \textbf{Automated Indicator Sharing (AIS), STIX/TAXII:} Tools and formats for sharing threat intelligence.
        \item \textbf{Predictive Analysis:} Forecasting future threats.
        \item \textbf{Threat Maps:} Visual representation of ongoing cyber attacks globally.
        \item \textbf{File/Code Repositories:} Places where software code is stored, which can sometimes contain vulnerabilities.
    \end{itemize}

    \subsubsection*{\color{purple}Research Sources}
    \begin{itemize}
        \item \textbf{Vendor Websites:} Companies that make software/hardware often provide updates or alerts.
        \item \textbf{Conferences:} Where experts discuss the latest in cybersecurity.
        \item \textbf{Academic Journals:} Peer-reviewed publications on new findings.
        \item \textbf{Request for Comments (RFC):} Official documentations and standards.
        \item \textbf{Local Industry Groups:} Local or regional groups focusing on security.
        \item \textbf{Social Media:} Real-time info, but needs verification.
        \item \textbf{Threat Feeds:} Live data streams about potential threats.
        \item \textbf{Adversary Tactics, Techniques, and Procedures (TTP):} Documented strategies used by attackers.
    \end{itemize}    
\end{fullwidth}
\newpage



\begin{fullwidth} 
    \subsection{Explain the security concerns associated with various types of vulnerabilities}
    \subsubsection*{\color{red}Actors and Threats}
    \begin{description}
        \item[Cloud-based vs. On-premises Vulnerabilities:] Cloud-based vulnerabilities: Relate to the weaknesses within cloud services and platforms that can be exploited by attackers, such as misconfigured cloud storage or inadequate identity and access management. On-premises vulnerabilities: Concern issues in your own physical environment (like a server room in your building), like outdated firewalls or servers with unpatched software.
        \item[Zero-day:] A Zero-day vulnerability refers to a software security flaw that is known to the software vendor but doesn’t have a patch in place to fix the vulnerability. It’s called “zero-day” because the developers have “zero days” to fix the problem that has just been exposed — and perhaps already exploited by hackers.
    \end{description}
    \subsubsection*{\color{red}Weak Configurations}
    \begin{itemize}
        \item \textbf{Open Permissions:} Allowing too much access to too many people/users.
        \item \textbf{Unsecure Root Accounts:} Not protecting high-level administrative accounts properly.
        \item \textbf{Errors:} Mistakes in coding or system setup.
        \item \textbf{Weak Encryption:} Not using strong methods to protect data.
        \item \textbf{Unsecure Protocols:} Using outdated or insecure communication protocols.
        \item \textbf{Default Settings:} Not changing the settings that the system or application came with.
        \item \textbf{Open Ports and Services:} Leaving too many openings for attackers to potentially exploit.
    \end{itemize}

    \subsubsection*{\color{red}Third-party Risks}
    \begin{itemize}
        \item \textbf{Vendor Management:} Not properly overseeing or managing the organizations you buy products or services from.
        \item \textbf{System Integration:} Problems that might arise when trying to get different systems to work together.
        \item \textbf{Lack of Vendor Support:} Vendors not providing sufficient help or updates for their products.
        \item \textbf{Supply Chain:} The process of creating and delivering a product, which can be disrupted or exploited at various stages.
        \item \textbf{Outsourced Code Development:} Getting external parties to write software for you, which might not be as secure.
        \item \textbf{Data Storage:} Where and how you store data, and the vulnerabilities there.
    \end{itemize}

    \begin{description}
        \item[Legacy Platforms:] Using outdated systems or software that no longer receive updates and therefore, might be full of vulnerabilities.] 
    \end{description}

    \subsubsection*{\color{red}Improper or Weak Patch Management}
    \begin{itemize}
        \item \textbf{Firmware:} The foundational software for hardware, often neglected in the patching process.
        \item \textbf{Operating System (OS):} The main software that runs a computer, which might be left outdated.
        \item \textbf{Applications:} Programs used for various purposes that might not be kept up-to-date with security patches.
    \end{itemize}

    \subsubsection*{\color{red}Impacts}
    \begin{itemize}
        \item \textbf{Firmware:} The foundational software for hardware, often neglected in the patching process.
        \item \textbf{Data Loss:} Losing data due to an incident.
        \item \textbf{Data Breaches:} Unauthorized access to data.
        \item \textbf{Data Exfiltration:} The unauthorized copying, transfer, or retrieval of data.
        \item \textbf{Identity Theft:} Unauthorized use of someone’s personal data.
        \item \textbf{Financial:} Monetary losses from an incident.
        \item \textbf{Reputation:} Damage to the organization’s standing.
        \item \textbf{Availability Loss:} Losing access to systems, data, or networks.
    \end{itemize}
\end{fullwidth}
\newpage
\begin{fullwidth} 
    \subsection{Summarize the techniques used in security assessments}
    \subsubsection*{\color{red}Threat Hunting}
    \begin{itemize}
        \item \textbf{Intelligence Fusion:} Combining various sources of information to generate actionable intelligence about threats.
        \item \textbf{Threat Feeds:} Streams of data related to potential threats, like IP addresses known to be malicious.
        \item \textbf{Advisories and Bulletins:} Alerts and notifications regarding new threats or vulnerabilities.
        \item \textbf{Maneuver:} Adapting to or moving against a threat to neutralize it.
    \end{itemize}
    \subsubsection*{\color{red}Vulnerability Scans}
    \begin{itemize}
        \item \textbf{False Positives:} Alerts on vulnerabilities that aren’t actually present (false alarms).
        \item \textbf{False Negatives:} Failing to detect an actual vulnerability (missing a real threat).
        \item \textbf{Log Reviews:} Analyzing logs to identify suspicious activity.
        \item \textbf{Credentialed vs. Non-credentialed:} Scans with login credentials vs. those without to see system vulnerabilities from different viewpoints.
        \item \textbf{Intrusive vs. Non-intrusive:} Scans that might impact system performance vs. those that don’t.
        \item \textbf{Application/Web Application/Network:} Scans targeting different elements: software applications, web platforms, or network infrastructure.
        \item \textbf{CVE/CVSS:} Standardized identifiers and scorings for vulnerabilities.
        \item \textbf{Configuration Review:} Checking system setups for vulnerabilities.
    \end{itemize}

    \subsubsection*{\color{red}Syslog/Security Information and Event Management (SIEM)}
    \begin{itemize}
        \item \textbf{Review Reports:} Analyzing compiled data and insights.
        \item \textbf{Packet Capture:} Collecting data packets transmitted over networks for analysis.
        \item \textbf{Data Inputs:} Different types of data fed into the SIEM for analysis.
        \item \textbf{User Behavior Analysis:} Studying how users interact with systems to identify anomalies.
        \item \textbf{Sentiment Analysis:} Utilizing data analysis to understand sentiments or attitudes expressed in source data.
        \item \textbf{Security Monitoring}: Continuously observing systems to detect and respond to security incidents.
        \item \textbf{Log Aggregation:} Collecting log data from different sources into a single location.
        \item \textbf{Log Collectors:} Systems or applications that gather log data.
    \end{itemize}

    \subsubsection*{\color{red}Security Orchestration, Automation, and Response (SOAR)}
    \begin{itemize}
        \item \textbf{Security Orchestration:} Coordinating and structuring how different security solutions work together.
        \item \textbf{Automation:} Utilizing technology to perform tasks without human intervention.
        \item \textbf{Response:} Actions taken to mitigate, prevent, or remediate security incidents.
    \end{itemize}

\end{fullwidth}
\newpage
\begin{fullwidth} 
    \subsection{Penetration Testing}
    \subsubsection*{\color{red}Threat Hunting}
    \begin{itemize}
        \item \textbf{Known/Unknown/Partially Known Environment:} Whether the tester has prior knowledge of the system, no knowledge, or some knowledge.
        \item \textbf{Rules of Engagement:} Agreements or boundaries defined before the test, such as what can be tested, when, and how.
        \item \textbf{Lateral Movement:} Moving from one network segment to another, often to gain access to a specific target.
        \item \textbf{Privilege Escalation:} Gaining higher-level privileges, such as moving from user-level to admin-level access.
        \item \textbf{Persistence:} Maintaining access to the system, even after the initial intrusion has been detected and supposedly “cleared”.
        \item \textbf{Cleanup:} Removing any traces and reversing any changes made during the test.
        \item \textbf{Bug Bounty:} Programs by organizations rewarding individuals for identifying and reporting vulnerabilities.
        \item \textbf{Pivoting:} Using a compromised system to target other systems.
    \end{itemize}
    \subsubsection*{\color{red}Passive and Active Reconnaissance}
    \begin{itemize}
        \item \textbf{Drones, War Flying, War Driving:} Using drones, aircraft, or vehicles to identify wireless networks.
        \item \textbf{Footprinting:} Gathering information about a target entity to find ways to infiltrate.
        \item \textbf{OSINT (Open-Source Intelligence):} Using publicly available sources to gather data about a target.
    \end{itemize}

    \subsubsection*{\color{red}Exercise Types}
    \begin{itemize}
        \item \textbf{Red Team:} Simulated attackers trying to compromise a system.
        \item \textbf{Blue Team:} The defenders who try to stop the red team.
        \item \textbf{White Team:} Neutral party, usually administrators, ensuring the test goes smoothly.
        \item \textbf{Purple Team:} A blend of red and blue teams, often where they work together to understand and remediate vulnerabilities.
    \end{itemize}

\end{fullwidth}
\newpage
\section{Architecture and Design}








\end{document}